\sectionkicker{Source-backed technical notes}
\subsection*{Open ModelsもAgent Stackへ — Qwen・GLM・MiniMax: Technical Notes}
本欄は記事本文で圧縮した一次資料上の情報を、比較・再検証しやすい形で整理したものである。時系列、確認済みの主張、留意点、一次資料URLを掲載する。
\medskip
\subsection*{Theme at a glance}
\addcontentsline{toc}{subsection}{Theme at a glance}
\begin{center}
\footnotesize
\begin{tabularx}{\linewidth}{@{}p{0.20\linewidth}p{0.10\linewidth}p{0.14\linewidth}X@{}}
\toprule
資料 & 位置づけ & 種別 & 時系列 \\
\midrule
Qwen3.5 / Qwen3-Coder-Next February Model Studio rollout & 主要資料 & モデル更新 & 2026-02-16 (REGIONAL\_モデル公開); 2026-02-20 (INTERNATIONAL\_モデル公開) \\
GLM-5 & 主要資料 & オープンウェイト & 2026-02-12 (オープンウェイト\_モデル公開) \\
MiniMax M2.5 + Forge agent-native RL & 主要資料 & モデル更新 & 2026-02-12 (モデル公開); 2026-02-13 (TECHNICAL\_Framework公開) \\
SGLang v0.5.9 & 補足資料 & Framework & 2026-02-24 (PROJECT\_RELEASE) \\
Transformers v5.1 / v5.2 model-family integrations & 補足資料 & Framework & 2026-02-05 (PROJECT\_RELEASE); 2026-02-16 (PROJECT\_RELEASE) \\
\bottomrule
\end{tabularx}
\end{center}
\normalsize

\begin{technicalnote}{Qwen3.5 / Qwen3-Coder-Next February Model Studio rollout}{主要資料}
\begingroup
% reader-facing Technical Notes break policy
\widowpenalty=10000
\clubpenalty=10000
\displaywidowpenalty=10000
\begin{tabularx}{\linewidth}{@{}>{\bfseries}p{0.22\linewidth}X@{}}
組織 & Alibaba Cloud / Qwen \\
種別 & モデル更新 \\
時系列 & 2026-02-16 (REGIONAL\_モデル公開); 2026-02-20 (INTERNATIONAL\_モデル公開) \\
\end{tabularx}
\smallskip
{\bfseries 一次資料から整理したtechnical points}
\begin{itemize}[leftmargin=1.5em,itemsep=0.35em]
\item \textbf{一次情報で確認できる事実}: Alibaba Cloud Model Studioのlifecycle記録では、qwen3.5-plusやQwen3.5-397B-A17Bを含むQwen3.5系列が2026年2月に利用可能となり、一次カタログにはtext・image・video入力とreasoning/coding/agent能力が記載されている。
\item \textbf{一次情報で確認できる事実}: international側のlifecycle記録では、Qwen3-Coder-Nextが2026年2月20日に、multi-turn tool interactionとrepository-level code understandingを対象とするopen-source code modelとして掲載されている。
\item \textbf{Vendor claim}: benchmark能力や効率の優位性に関する主張は、独立再現されない限りvendor claimとして扱う。
\end{itemize}
\begin{minipage}{\linewidth}
% reader-facing Technical Notes generic boundary/source tail group
{\bfseries 読む際の境界}
\begin{itemize}[leftmargin=1.5em,itemsep=0.35em]
\item \textbf{分析上の留意点}: Model Studioのavailabilityと日付はChina/internationalで異なるため、各lifecycle entryの地域範囲を保持する必要がある。
\item \textbf{分析上の留意点}: lifecycle pageは継続更新され後日のglobal rolloutを含み得るため、2月のchronologyは現在のavailabilityではなく日付付きentryを基準にする。
\end{itemize}
\begin{samepage}
{\bfseries 一次資料}
\begingroup\sloppy
\Urlmuskip=0mu plus 2mu\relax
\begin{itemize}[leftmargin=1.5em,itemsep=0.25em]
\item {\scriptsize\url{https://www.alibabacloud.com/help/en/model-studio/newly-released-models}}
\end{itemize}
\endgroup
\end{samepage}
\end{minipage}
\endgroup
\end{technicalnote}
\medskip

\begin{technicalnote}{GLM-5}{主要資料}
\begingroup
% reader-facing Technical Notes break policy
\widowpenalty=10000
\clubpenalty=10000
\displaywidowpenalty=10000
\begin{tabularx}{\linewidth}{@{}>{\bfseries}p{0.22\linewidth}X@{}}
組織 & Z.ai \\
種別 & オープンウェイト \\
時系列 & 2026-02-12 (オープンウェイト\_モデル公開) \\
\end{tabularx}
\smallskip
{\bfseries 一次資料から整理したtechnical points}
\begin{itemize}[leftmargin=1.5em,itemsep=0.35em]
\item \textbf{一次情報で確認できる事実}: Z.aiは2026年2月12日にGLM-5を公開し、complex systems engineeringとlong-horizon agentic taskを対象とするmodelとして説明した。
\item \textbf{一次情報で確認できる事実}: first-party releaseでは、総744B parameters・40B active、28.5T pre-training tokens、DeepSeek Sparse Attentionの統合、slimeと呼ばれるasynchronous RL infrastructureが記載されている。
\item \textbf{一次情報で確認できる事実}: 公開weightsはMIT licenseと説明され、vLLMやSGLangを含むframeworkによるdeployment pathが示されている。
\item \textbf{Vendor claim}: benchmark rankingや、GLM-5がfrontier proprietary modelとの差を縮めたとする主張はvendor claimである。
\end{itemize}
\begin{minipage}{\linewidth}
% reader-facing Technical Notes generic boundary/source tail group
{\bfseries 読む際の境界}
\begin{itemize}[leftmargin=1.5em,itemsep=0.35em]
\item \textbf{分析上の留意点}: architecture/trainingの数値はfirst-party technical disclosureであり、独立監査されたものではない。
\item \textbf{分析上の留意点}: benchmark上の優位性に関する記述はvendor報告のままである。
\item \textbf{分析上の留意点}: modelのreleaseとlicenseは、主張されるperformanceとは独立して検証できる。
\end{itemize}
\begin{samepage}
{\bfseries 一次資料}
\begingroup\sloppy
\Urlmuskip=0mu plus 2mu\relax
\begin{itemize}[leftmargin=1.5em,itemsep=0.25em]
\item {\scriptsize\url{https://docs.z.ai/release-notes/new-released}}
\end{itemize}
\endgroup
\end{samepage}
\end{minipage}
\endgroup
\end{technicalnote}
\medskip

\begin{technicalnote}{MiniMax M2.5 + Forge agent-native RL}{主要資料}
\begingroup
% reader-facing Technical Notes break policy
\widowpenalty=10000
\clubpenalty=10000
\displaywidowpenalty=10000
\begin{tabularx}{\linewidth}{@{}>{\bfseries}p{0.22\linewidth}X@{}}
組織 & MiniMax \\
種別 & モデル更新 \\
時系列 & 2026-02-12 (モデル公開); 2026-02-13 (TECHNICAL\_Framework公開) \\
\end{tabularx}
\smallskip
{\bfseries 一次資料から整理したtechnical points}
\begin{itemize}[leftmargin=1.5em,itemsep=0.35em]
\item \textbf{一次情報で確認できる事実}: MiniMaxは2026年2月12日にM2.5を発表し、coding、agentic tool use/search、office/workflow taskを重点領域として示した。
\item \textbf{一次情報で確認できる事実}: MiniMaxは2026年2月13日にForgeを、scaffold/environment abstractionと大規模agent training向けasynchronous schedulingを備えたagent-native reinforcement-learning frameworkとして紹介した。
\item \textbf{Vendor claim}: M2.5のbenchmark score・speedup、およびForgeのscale/throughput値はvendorによる報告である。
\end{itemize}
\begin{minipage}{\linewidth}
% reader-facing Technical Notes generic boundary/source tail group
{\bfseries 読む際の境界}
\begin{itemize}[leftmargin=1.5em,itemsep=0.35em]
\item \textbf{分析上の留意点}: M2.5のbenchmarkおよびrelative-speedの数値は独立再現していない。
\item \textbf{分析上の留意点}: Forgeのscaleおよびtraining-throughputに関する主張はfirst-party reportである。
\item \textbf{分析上の留意点}: Forge自体が広く利用可能なproductであるかのように見せず、public model releaseとtraining-systemの話を区別する必要がある。
\end{itemize}
\begin{samepage}
{\bfseries 一次資料}
\begingroup\sloppy
\Urlmuskip=0mu plus 2mu\relax
\begin{itemize}[leftmargin=1.5em,itemsep=0.25em]
\item {\scriptsize\url{https://www.minimax.io/news}}
\end{itemize}
\endgroup
\end{samepage}
\end{minipage}
\endgroup
\end{technicalnote}
\medskip

\begin{technicalnote}{SGLang v0.5.9}{補足資料}
\begingroup
% reader-facing Technical Notes break policy
\widowpenalty=10000
\clubpenalty=10000
\displaywidowpenalty=10000
\begin{tabularx}{\linewidth}{@{}>{\bfseries}p{0.22\linewidth}X@{}}
組織 & SGLang Project \\
種別 & Framework \\
時系列 & 2026-02-24 (PROJECT\_RELEASE) \\
\end{tabularx}
\smallskip
{\bfseries 一次資料から整理したtechnical points}
\begin{itemize}[leftmargin=1.5em,itemsep=0.35em]
\item \textbf{一次情報で確認できる事実}: SGLang v0.5.9は2026年2月24日に公開され、Kimi-K2.5、GLM-5、Qwen3.5、MiniMax M2.5、および複数のmultimodal/diffusionモデルファミリへの対応が追加された。
\item \textbf{一次情報で確認できる事実}: このreleaseでは、Anthropic互換API endpointに加え、inference、MoE、multimodal、diffusion runtimeに関する大幅な変更が追加された。
\item \textbf{Project claim}: overlapped LoRA loadingでTTFTを約78\%低減、Blackwell上でDeepSeek V3.2 NSAを3〜5倍高速化したとする数値は、独立測定ではなくproject側の報告である。
\end{itemize}
\begin{minipage}{\linewidth}
% reader-facing Technical Notes generic boundary/source tail group
{\bfseries 読む際の境界}
\begin{itemize}[leftmargin=1.5em,itemsep=0.35em]
\item \textbf{分析上の留意点}: 性能値はprojectのrelease noteに基づくもので、本調査では独立に再現していない。
\item \textbf{分析上の留意点}: support追加は実装が利用可能になったことを示すが、model family間の品質同等性やbenchmarkの妥当性を証明するものではない。
\end{itemize}
\begin{samepage}
{\bfseries 一次資料}
\begingroup\sloppy
\Urlmuskip=0mu plus 2mu\relax
\begin{itemize}[leftmargin=1.5em,itemsep=0.25em]
\item {\scriptsize\url{https://github.com/sgl-project/sglang/releases/tag/v0.5.9}}
\end{itemize}
\endgroup
\end{samepage}
\end{minipage}
\endgroup
\end{technicalnote}
\medskip

\begin{technicalnote}{Transformers v5.1 / v5.2 model-family integrations}{補足資料}
\begingroup
% reader-facing Technical Notes break policy
\widowpenalty=10000
\clubpenalty=10000
\displaywidowpenalty=10000
\begin{tabularx}{\linewidth}{@{}>{\bfseries}p{0.22\linewidth}X@{}}
組織 & Hugging Face \\
種別 & Framework \\
時系列 & 2026-02-05 (PROJECT\_RELEASE); 2026-02-16 (PROJECT\_RELEASE) \\
\end{tabularx}
\smallskip
{\bfseries 一次資料から整理したtechnical points}
\begin{itemize}[leftmargin=1.5em,itemsep=0.35em]
\item \textbf{一次情報で確認できる事実}: Transformers v5.2.0は2026年2月16日に公開され、VoxtralRealtime、GLM-5、Qwen3.5/Qwen3.5 MoE、VibeVoice Acoustic Tokenizerへの明示的なsupportが追加された。
\item \textbf{Project claim}: VoxtralRealtime implementationでは、audio encoderとMistral-based language-model decoderにtime conditioningとcausal-convolution cacheを組み合わせたstreaming ASRが説明されている。
\item \textbf{一次情報で確認できる事実}: Transformers v5.1.0は2月前半に公開され、GLM-OCRなど追加のmodel-family integrationを含んでいた。
\item \textbf{分析上の留意点}: Transformers release note内に引用されたvendorのperformance記述は、Hugging Faceによる独立検証ではない。
\end{itemize}
\begin{minipage}{\linewidth}
% reader-facing Technical Notes generic boundary/source tail group
{\bfseries 読む際の境界}
\begin{itemize}[leftmargin=1.5em,itemsep=0.35em]
\item \textbf{分析上の留意点}: framework integrationはTransformersへimplementation/supportが入ったことを示すが、model vendorのbenchmark claimを独立検証するものではない。
\item \textbf{分析上の留意点}: v5.1.0はsupporting contextであり、2月のintegration milestoneとしてはv5.2.0の方が強い。
\item \textbf{分析上の留意点}: v5.2 releaseにはbreaking interface changeも含まれるため、model supportを単純なcompatibility listだけに還元すべきではない。
\end{itemize}
\begin{samepage}
{\bfseries 一次資料}
\begingroup\sloppy
\Urlmuskip=0mu plus 2mu\relax
\begin{itemize}[leftmargin=1.5em,itemsep=0.25em]
\item {\scriptsize\url{https://github.com/huggingface/transformers/releases/tag/v5.1.0}}
\item {\scriptsize\url{https://github.com/huggingface/transformers/releases/tag/v5.2.0}}
\end{itemize}
\endgroup
\end{samepage}
\end{minipage}
\endgroup
\end{technicalnote}
\medskip
